\begingroup
\color{black}

\chapter{Discussion and Conclusions}
\label{ch:discussion-conclusion}

This thesis asks whether graph-based decentralized policies can control one
grid and remain useful on a larger one. The work addresses this question in
four stages: stable learning, sparse intervention, graph design, and
cross-grid transfer.

\section{Research questions revisited}

\noindent\textbf{1. Which changes stabilize decentralized MAPPO?}
The tuned flat actor reaches high survival earlier and more consistently than
the initial implementation. Its final mean survival is about $93\%$ on
\texttt{bus14}. Reward normalization, the critic update, and the optimization
settings have a clearer effect than simply increasing the model size. An
initial bias toward do nothing is not consistently helpful: a strong bias
delays learning, while a weaker bias remains unstable. The zero-bias model is
therefore retained as a stable reference, even though it does not have the
highest endpoint in every comparison.

\medskip
\noindent\textbf{2. Can unnecessary interventions be reduced?}
Yes, at least on the source grid. The local-safe adaptive intervention budget
reaches about $98.4\%$ mean survival while selecting action 0 in about $98\%$
of decisions. The baseline reaches $93.67\%$ survival and selects action 0 in
only $43.4\%$ of decisions. The agents can therefore learn to intervene less
often without losing performance. A fixed penalty also gives good results,
whereas the learned gate does not improve both survival and intervention
frequency at the same time.

\medskip
\noindent\textbf{3. Which graph design provides a useful local actor?}
Several graph actors exceed $99\%$ mean survival on \texttt{bus14}. For the
fixed-head encoder branch, the source-grid screens retain GINE message passing,
two layers of width 128, and an energized-node maximum graph readout. Several
busbar and disaggregated representations are then kept because their small
source-grid differences may become more important after transfer. The
candidate-action head also works well: six of its eight cells exceed $98.2\%$
and three exceed $99\%$. More importantly, its shared scorer can rank action
lists of different sizes, which makes complete-policy transfer possible.

\medskip
\noindent\textbf{4. Do graph-policy weights transfer to WCCI?}
The complete shared encoder and candidate scorer can run on WCCI without any
change to their parameter dimensions. This confirms structural transfer, but
the transferred behavior remains limited. On the 24 difficult chronics, do
nothing reaches $8.34\%$ mean survival. With the local-$\rho=0.95$ heuristic,
the median zero-shot actor reaches $21.9\%$ at $k=32$, while the best evaluated
actor reaches $25.8\%$. Performance decreases as the candidate list grows even
though the greedy reference improves. The main difficulty is therefore not
the lack of useful actions, but ranking them correctly. In the matched $k=64$
comparison, fine-tuning raises median difficult-chronic survival to $18.1\%$,
compared with $9.6\%$ from scratch and $5.7\%$ zero-shot. Source initialization
is useful, but target adaptation is still needed.
\\
\\
\textcolor{red}{This answer will be completed after Chapter~\ref{ch:encoder-transfer-bridge}
establishes whether a frozen source encoder outperforms the random-frozen
control when the WCCI action heads are relearned.}

\section{What the experiments show together}

High source-grid survival does not guarantee successful transfer. The screens
in Chapter~\ref{ch:experiments} show that the graph encoders and candidate
scorer can learn the \texttt{bus14} task. The WCCI experiments then test two
different forms of reuse. Encoder-only transfer asks whether the learned
representation remains useful after replacing the grid-specific head.
Complete-policy transfer also reuses the action-selection mechanism.
\\
\\
Together, the experiments give three practical lessons. First, compatible
parameter dimensions do not guarantee compatible behavior. The input
convention, candidate actions, and intervention frequency all matter. Second,
a smaller candidate set can help a learned policy even when a larger set
contains better greedy actions, because the policy must still rank those
actions correctly. Third, the local heuristic is policy-dependent: it improves
most zero-shot actors but can hurt policies trained directly on WCCI. Overall
survival should therefore be read together with difficult-chronic survival,
since many of the easy chronics are already completed by doing nothing.

\section{Limitations and next steps}

The main architecture screens use one training seed, and the controlled WCCI
comparison uses 50 chronics from a single source-target grid pair. Maintenance
is disabled on WCCI and the policies use reduced action spaces. Scratch and
fine-tuned actors also receive different training budgets. Their comparison
therefore supports the observed performance ordering, but not a strict claim
about sample efficiency. Finally, the one-step greedy controller is a useful
reference, not an upper bound on long-term control.
\\
\\
The next experiments should focus on the finalists rather than repeat the full
architecture search. They should use additional seeds, a larger untouched WCCI
cohort, and equal budgets for scratch training and fine-tuning. Testing with
maintenance and on a second target grid would show whether the conclusions are
more general. Better candidate ranking is also a priority, because the largest
remaining gap lies between the learned policies and the useful actions already
available to them.

\section{Conclusion}

Graph and candidate-action architectures remove the dimensional barrier that
prevents a fixed-width decentralized actor from being reused across grids.
However, structural compatibility alone does not guarantee reliable control
on a new network. In the \texttt{bus14}-to-WCCI study, useful transfer also
depends on compatible inputs, a manageable action set, candidate-aware
scoring, policy-specific intervention control, and target adaptation. The main
conclusion is therefore not that one graph architecture solves transfer. It is
that a transferable policy can retain meaningful zero-shot performance and
provide a better starting point for adaptation than learning the same control
problem again from random weights.

\endgroup
